% !TEX encoding = UTF-8 Unicode
\documentclass[12pt,a4paper,oneside,brazil,chapter=TITLE]{article}
\usepackage[T1]{fontenc}
\usepackage[utf8]{inputenc}
\usepackage[brazil]{babel}
\usepackage{times}
\renewcommand{\ttdefault}{cmtt}
\usepackage[a4paper,left=3cm,top=3cm,right=2cm,bottom=2cm,headheight=14pt,headsep=0.5cm]{geometry}
\usepackage{setspace}
\onehalfspacing
\usepackage{indentfirst}
\setlength{\parindent}{1.25cm}
\setlength{\parskip}{0pt}
\usepackage{titlesec}
\titleformat{\section}{\normalfont\normalsize\bfseries\MakeUppercase}{\thesection\quad}{0em}{}
\titleformat{\subsection}{\normalfont\normalsize\bfseries}{\thesubsection\quad}{0em}{}
\titleformat{\subsubsection}{\normalfont\normalsize\bfseries\itshape}{\thesubsubsection\quad}{0em}{}
\titlespacing*{\section}{0pt}{18pt}{12pt}
\titlespacing*{\subsection}{0pt}{18pt}{6pt}
\titlespacing*{\subsubsection}{0pt}{12pt}{6pt}
\newcommand{\secbreak}{\clearpage}
\let\oldsection\section
\renewcommand{\section}{\secbreak\oldsection}
\usepackage{fancyhdr}
\pagestyle{fancy}
\fancyhf{}
\fancyhead[R]{\fontsize{10}{12}\selectfont\thepage}
\renewcommand{\headrulewidth}{0pt}
\renewcommand{\footrulewidth}{0pt}
\fancypagestyle{plain}{\fancyhf{}\fancyhead[R]{\fontsize{10}{12}\selectfont\thepage}\renewcommand{\headrulewidth}{0pt}}
\usepackage{longtable}
\usepackage{booktabs}
\usepackage{array}
\usepackage{multirow}
\usepackage{tabularx}
\renewcommand{\arraystretch}{1.3}
\usepackage{caption}
\DeclareCaptionFormat{abnt}{\fontsize{10}{12}\selectfont#1#2#3}
\DeclareCaptionLabelSeparator{emdashsep}{ --- }
\captionsetup{format=abnt,justification=centering,singlelinecheck=false,labelsep=emdashsep,position=above,skip=3pt}
\newcommand{\fonte}[1]{\begin{flushleft}\fontsize{10}{12}\selectfont Fonte: #1\end{flushleft}}
\usepackage{enumitem}
\setlist{nosep, leftmargin=*}
\providecommand{\tightlist}{\setlength{\itemsep}{0pt}\setlength{\parskip}{0pt}}
\usepackage{tocloft}
\renewcommand{\contentsname}{SUMÁRIO}
\newlistof{quadros}{loq}{LISTA DE QUADROS}
\setlength{\cftquadrosindent}{0pt}
\setlength{\cftquadrosnumwidth}{0pt}
\makeatletter
\renewcommand{\listofquadros}{\section*{LISTA DE QUADROS}\@starttoc{loq}}
\makeatother
\newcommand{\quadrotitle}[2]{\vspace{1em}\phantomsection\label{quadro#1}\addcontentsline{loq}{quadros}{Quadro #1 --- #2}\noindent\textbf{Quadro #1 --- #2}\par\vspace{0.3em}}
\newcommand{\quadrofonte}[1]{\par\noindent\textit{Fonte: #1}\vspace{1em}\par}
\PassOptionsToPackage{hyphens}{url}
\usepackage[hidelinks,colorlinks=false,breaklinks=true]{hyperref}
\usepackage{quoting}
\quotingsetup{font=small,indentfirst=false,vskip=\baselineskip}
\usepackage{graphicx}
\usepackage{float}
\date{}

% ============================================================
% DOCUMENTO
% ============================================================
\begin{document}

% ----------------------------------------------------------
% CAPA
% ----------------------------------------------------------
\begin{titlepage}
  \thispagestyle{empty}
  \begin{center}
    {\normalsize\bfseries
      Instituto Federal de Educação, Ciência e Tecnologia de Brasília --- IFB\\[0.3em]
      \textit{Campus} Samambaia\\[0.3em]
      Curso Superior de Tecnologia em Design de Produto
    }
    \vfill
    {\large\bfseries RELATÓRIO DE VISITA TÉCNICA --- GRAVIA INDÚSTRIA DE PERFILADOS DE AÇO}
    \vfill
    {\normalsize Brasília/DF\\2026}
  \end{center}
\end{titlepage}
\setcounter{page}{2}

% ----------------------------------------------------------
% FOLHA DE ROSTO
% ----------------------------------------------------------
\begin{titlepage}
  \thispagestyle{empty}
  \begin{center}
    {\normalsize\bfseries VICTOR HUGO DA SILVA COSTA}
    \vfill
    {\large\bfseries RELATÓRIO DE VISITA TÉCNICA --- GRAVIA INDÚSTRIA DE PERFILADOS DE AÇO}
    \vfill
    \begin{flushright}
      \begin{minipage}{0.5\textwidth}
        \normalsize
        Relatório de visita técnica apresentado ao Curso Superior de Tecnologia em Design de Produto do Instituto Federal de Educação, Ciência e Tecnologia de Brasília --- IFB, \textit{Campus} Samambaia, como requisito parcial para a aprovação na disciplina de Materiais e Processos de Fabricação I.
        \vspace{1em}

        Professora: Prof.ª Keila Sanches
      \end{minipage}
    \end{flushright}
    \vfill
    {\normalsize Brasília/DF --- 2026}
  \end{center}
\end{titlepage}

% ----------------------------------------------------------
% RESUMO
% ----------------------------------------------------------
\clearpage
\thispagestyle{plain}
\begin{center}
  {\normalsize\bfseries RESUMO}
\end{center}
\vspace{12pt}

\noindent Este relatório descreve a visita técnica realizada em 02 de junho de 2026 à unidade industrial da Gravia Indústria de Perfilados de Aço, em Taguatinga, Distrito Federal. O objetivo foi observar os processos de beneficiamento do aço executados pela empresa, compreender a organização dos setores produtivos e identificar os principais equipamentos empregados na fabricação de componentes para a construção civil e aplicações industriais. A planta é dividida em galpões especializados: o Galpão 1 concentra operações de corte e conformação de chapas; o Galpão 3 é dedicado à perfilação contínua para sistemas de drywall; e os Galpões 4, 5 e 6 atendem ao processamento de chapas espessas. Os equipamentos registrados incluem prensas dobradeiras hidráulicas, corte a laser, guilhotina hidráulica, corte a plasma, perfiladeiras e pontes rolantes. A gestão da produção adota Lean Six Sigma e Gestão por Processos e Trabalho Padronizado (GPTW).

\vspace{12pt}
\noindent\textbf{Palavras-chave:} Beneficiamento do aço; Processos de fabricação; Corte a laser; Conformação por dobra; Perfilação.

% ----------------------------------------------------------
% LISTA DE QUADROS
% ----------------------------------------------------------
\clearpage
\listofquadros

% ----------------------------------------------------------
% SUMÁRIO
% ----------------------------------------------------------
\clearpage
\tableofcontents

% ============================================================
% CORPO DO TEXTO
% ============================================================

% ----------------------------------------------------------
% SEÇÃO 1 — INTRODUÇÃO
% ----------------------------------------------------------
\section{Introdução}

A visita técnica à unidade industrial da Gravia Indústria de Perfilados de Aço foi realizada em 02 de junho de 2026, com o propósito de observar os processos de beneficiamento do aço em operação, compreender a organização funcional da planta e registrar os equipamentos empregados nas diferentes etapas de fabricação. Com mais de seis décadas de atuação no mercado e unidade instalada em Taguatinga, Distrito Federal, a empresa se especializou na transformação do aço para os setores da construção civil, da indústria e do agronegócio.

% ----------------------------------------------------------
% SEÇÃO 2 — ORGANIZAÇÃO DA PLANTA E SISTEMA DE GESTÃO
% ----------------------------------------------------------
\section{Organização da Planta e Sistema de Gestão}

A planta é dividida em galpões com funções bem delimitadas, o que permite controle mais fino do fluxo de materiais e especialização das equipes por setor. Essa compartimentação ficou evidente já na entrada da fábrica, onde placas indicativas e codificação visual delimitavam com clareza as áreas de circulação e os postos de trabalho.

Na gestão da produção, a empresa combina o Lean Six Sigma, cujo foco é eliminar desperdícios e reduzir a variabilidade dos processos por meios estatísticos, com a Gestão por Processos e Trabalho Padronizado (GPTW), que documenta rotinas operacionais para cada etapa produtiva. O resultado prático dessas metodologias era visível no chão de fábrica: materiais identificados, etapas registradas e espaços organizados sem excesso.

A empresa mantém também o Setor de Tubos (SIA), unidade voltada à fabricação e ao processamento de tubos metálicos para clientes internos e externos.

% ----------------------------------------------------------
% SEÇÃO 3 — GALPÃO 1
% ----------------------------------------------------------
\section{Galpão 1 --- Corte e Conformação de Chapas}

O Galpão 1 é o principal centro de transformação de chapas metálicas da fábrica, com a maior concentração e diversidade de equipamentos observados na visita. Foi o que ficou mais na memória do percurso: o ruído contínuo das prensas dobradeiras Newton em operação, a cadência das chapas sendo alimentadas pelas desbobinadeiras e o volume de peças semi-acabadas empilhadas ao longo do galpão.

\subsection{Fluxo Produtivo}

O fluxo produtivo do Galpão 1 segue a seguinte sequência operacional:

\begin{enumerate}
\def\labelenumi{\arabic{enumi}.}
\tightlist
\item Recebimento das bobinas de aço;
\item Desbobinamento e alimentação das linhas de produção;
\item Corte das bobinas em chapas padronizadas de 3 m e 6 m;
\item Processos de conformação e fabricação das peças finais.
\end{enumerate}

\subsection{Equipamentos}

\quadrotitle{1}{Equipamentos do Galpão 1 --- corte e conformação de chapas}
{\setlength{\tabcolsep}{4pt}
\begin{longtable}{@{}
  >{\raggedright\arraybackslash}p{0.24\linewidth}
  >{\raggedright\arraybackslash}p{0.20\linewidth}
  >{\raggedright\arraybackslash}p{0.44\linewidth}
  >{\centering\arraybackslash}p{0.06\linewidth}@{}}
\toprule
\textbf{Equipamento} & \textbf{Modelo} & \textbf{Características} & \textbf{Qtd.} \\
\midrule
\endfirsthead
\toprule
\textbf{Equipamento} & \textbf{Modelo} & \textbf{Características} & \textbf{Qtd.} \\
\midrule
\endhead
\bottomrule
\endlastfoot
Prensa dobradeira hidráulica & Newton PSH-35030 & Comprimento útil aprox. 3 m & 4 \\
Prensa dobradeira hidráulica & Newton PSH-7030 & Comprimento útil aprox. 6 m & 6 \\
Corte a laser & Newton LN 3015F & Alta precisão dimensional; programação por software; aplicação em geometrias complexas & 1 \\
Prensa hidráulica sincronizada & Newton PSH-25030 & Controle de precisão em operações de dobra específicas & 1 \\
Guilhotina hidráulica & Newton GHN 301911 (3019) & Aplicação em chapas de maiores espessuras & 1 \\
Ponte rolante & HTB 1787 & Movimentação de chapas e componentes de grande porte & 1 \\
Desbobinadeira & --- & Alimentação das linhas; conversão de bobinas em chapas & 1 \\
Torno mecânico & --- & Fabricação e manutenção de peças e dispositivos internos & 1 \\
\end{longtable}}
\quadrofonte{elaborado pelo autor com base em observação \textit{in loco} (2026).}

O equipamento de maior interesse técnico do setor é o corte a laser Newton LN 3015F. Por ser programado via software específico, permite a execução de geometrias complexas com alta repetibilidade dimensional, algo inviável por guilhotina ou plasma nos mesmos parâmetros. A presença de um torno mecânico no galpão também chama atenção: indica que a empresa fabrica e mantém seus próprios dispositivos internos, o que aponta para um grau relevante de verticalização da manutenção.

% ----------------------------------------------------------
% SEÇÃO 4 — GALPÃO 3
% ----------------------------------------------------------
\section{Galpão 3 --- Perfilação e Produtos para Drywall}

O Galpão 3 é dedicado à fabricação de perfis metálicos para sistemas construtivos em drywall, com produção em série a partir de bobinas de aço processadas em regime contínuo.

\subsection{Produtos Fabricados}

Os produtos deste setor incluem:

\begin{itemize}
\tightlist
\item Trilhos para drywall;
\item Montantes metálicos;
\item Perfis metálicos diversos;
\item Tubos de aproximadamente 2'' de diâmetro.
\end{itemize}

\subsection{Equipamentos}

\quadrotitle{2}{Equipamentos do Galpão 3 --- perfilação e produtos para drywall}
{\setlength{\tabcolsep}{4pt}
\begin{longtable}{@{}
  >{\raggedright\arraybackslash}p{0.30\linewidth}
  >{\raggedright\arraybackslash}p{0.64\linewidth}@{}}
\toprule
\textbf{Equipamento} & \textbf{Características} \\
\midrule
\endfirsthead
\toprule
\textbf{Equipamento} & \textbf{Características} \\
\midrule
\endhead
\bottomrule
\endlastfoot
Tesoura rotativa & Processamento contínuo de bobinas metálicas \\
Perfiladeira & Velocidade aprox. 32 m/min; largura máx. 1.200 mm; espessuras entre 0,9 mm e 5 mm \\
\end{longtable}}
\quadrofonte{elaborado pelo autor com base em observação \textit{in loco} (2026).}

A perfiladeira opera em regime contínuo. A velocidade de processamento de cerca de 32 m/min, combinada com a faixa de espessuras entre 0,9 mm e 5 mm, confere ao equipamento boa versatilidade dentro dos limites da linha. Do ponto de vista de design de produto, é interessante observar que a padronização dimensional obtida por esse processo é condição necessária para a compatibilidade dos sistemas construtivos em drywall, onde a intercambialidade entre perfis de diferentes fornecedores depende de tolerâncias apertadas.

% ----------------------------------------------------------
% SEÇÃO 5 — GALPÕES 4, 5 E 6
% ----------------------------------------------------------
\section{Galpões 4, 5 e 6 --- Processamento de Chapas Espessas}

Esses galpões concentram a fabricação de componentes estruturais para construção civil e aplicações industriais, com chapas de espessuras superiores às trabalhadas no Galpão 1.

\subsection{Equipamentos}

\quadrotitle{3}{Equipamentos dos Galpões 4, 5 e 6 --- processamento de chapas espessas}
{\setlength{\tabcolsep}{4pt}
\begin{longtable}{@{}
  >{\raggedright\arraybackslash}p{0.22\linewidth}
  >{\raggedright\arraybackslash}p{0.24\linewidth}
  >{\raggedright\arraybackslash}p{0.48\linewidth}@{}}
\toprule
\textbf{Equipamento} & \textbf{Modelo} & \textbf{Características} \\
\midrule
\endfirsthead
\toprule
\textbf{Equipamento} & \textbf{Modelo} & \textbf{Características} \\
\midrule
\endhead
\bottomrule
\endlastfoot
Corte a plasma & Oxipira Série Master & Capacidade de corte: 0,5 mm a 50 mm; aplicação em componentes estruturais; maior produtividade para chapas espessas; menor precisão dimensional que o laser \\
Prensa dobradeira hidráulica & Newton PSH-55060 & Comprimento útil até 6 m; aplicação em componentes estruturais de grandes dimensões \\
\end{longtable}}
\quadrofonte{elaborado pelo autor com base em observação \textit{in loco} (2026).}

A escolha pelo corte a plasma neste setor, em detrimento do laser utilizado no Galpão 1, reflete uma decisão de processo coerente: para chapas espessas, o plasma entrega produtividade maior a um custo operacional inferior, aceitando-se a perda de tolerância dimensional que, para componentes estruturais, normalmente não é fator crítico. Aí fica claro como a seleção do processo de fabricação não é questão de preferência tecnológica, mas de adequação entre as características do material, a geometria do produto e as exigências dimensionais da aplicação.

% ----------------------------------------------------------
% SEÇÃO 6 — CONSIDERAÇÕES FINAIS
% ----------------------------------------------------------
\section{Considerações Finais}

A visita à Gravia permitiu observar, em operação real, o fluxo de beneficiamento do aço desde o recebimento das bobinas até componentes acabados prontos para expedição. A experiência evidenciou algo que a teoria por si só não transmite bem: a lógica de por que cada processo está onde está dentro da planta.

A coexistência de corte a laser e corte a plasma no mesmo ambiente produtivo, por exemplo, ilustra de forma concreta os trade-offs entre precisão e produtividade que permeiam as decisões de engenharia de processo. Isso, junto com o volume de prensas dobradeiras de diferentes portes, mostra que a planta foi configurada menos em torno de uma tecnologia única e mais em função da diversidade de produtos que a empresa precisa atender.

A organização por galpões especializados, combinada com as metodologias Lean Six Sigma e GPTW em operação visível, compõe um modelo produtivo que parece ter amadurecido ao longo das décadas de atuação da empresa, e não apenas adotado formalmente.

% ----------------------------------------------------------
% REFERÊNCIAS
% ----------------------------------------------------------
\section{REFERÊNCIAS}\label{referencias}

\noindent GRAVIA INDÚSTRIA DE PERFILADOS DE AÇO. \textbf{Visita técnica à unidade industrial de Taguatinga/DF.} [Observação direta]. Taguatinga, DF, 02 jun. 2026.

\end{document}
